\documentclass[../DoAn.tex]{subfiles}
\begin{document}

Chương 1 đã trình bày bối cảnh, mục tiêu và định hướng chung của đề tài. Tiếp nối phần mở đầu đó, Chương 2 tập trung làm rõ nền tảng lý thuyết cần thiết cho bài toán vận hành MLS trên mạng ngang hàng. Trọng tâm của chương không phải là liệt kê thật nhiều kiến thức rời rạc, mà là chỉ ra mối liên hệ giữa ba lớp vấn đề: lõi mật mã của MLS, đặc tính bất đồng bộ của mạng P2P và yêu cầu duy trì một trạng thái nhóm nhất quán.

Nội dung chương được tổ chức thành bốn phần. Trước hết, chương phân tích ngữ cảnh bài toán để làm rõ vì sao khi đưa MLS từ mô hình có máy chủ trung tâm sang mô hình ngang hàng, khó khăn lớn nhất không nằm ở việc truyền tin mà nằm ở việc giữ cho các thành viên cùng đi trên một chuỗi trạng thái mật mã hợp lệ. Tiếp theo, chương khảo sát các hướng nghiên cứu liên quan nhằm chỉ ra ưu điểm, giới hạn và khoảng trống còn lại. Sau đó, chương trình bày các kiến thức nền tảng quan trọng nhất của MLS và TreeKEM. Cuối cùng, chương giới thiệu ngắn gọn những đặc tính của mạng ngang hàng, cơ chế phát tán thông điệp và công cụ hỗ trợ sắp xếp thông điệp ứng dụng trong hệ phân tán.

\section{Ngữ cảnh của bài toán: nhất quán trạng thái mật mã trong hệ phân tán}
\label{sec:ngu-canh-bai-toan}

Trong hệ phân tán nói chung, bài toán nhất quán trạng thái nhằm bảo đảm nhiều nút cùng hiểu và cùng cập nhật dữ liệu theo một quy tắc thống nhất. Đối với hệ thống nhắn tin nhóm được mã hóa đầu cuối, trạng thái cần được giữ nhất quán không chỉ là dữ liệu ứng dụng, mà còn là trạng thái mật mã dùng để sinh khóa, xác thực thành viên và ràng buộc lịch sử hội thoại \cite{rfc9420,rfc9750}. Vì vậy, nếu các nút không còn cùng nhìn thấy một trạng thái mật mã chung, khả năng giao tiếp của nhóm sẽ bị ảnh hưởng trực tiếp.

Trong MLS, trạng thái nhóm tiến hóa theo các kỷ nguyên, gọi là epoch. Có thể hiểu mỗi epoch là một phiên bản cụ thể của nhóm, trong đó tập thành viên hợp lệ cùng chia sẻ một bộ bí mật tương ứng. Khi nhóm thêm thành viên, loại bỏ thành viên hoặc cập nhật khóa, trạng thái này phải chuyển sang epoch mới. Nếu tất cả thành viên cùng áp dụng các thay đổi theo cùng một thứ tự, nhóm tiếp tục hoạt động bình thường. Ngược lại, nếu các thay đổi được áp dụng theo những thứ tự khác nhau, các thành viên có thể dần tách khỏi nhau về mặt mật mã \cite{rfc9420,alwen2023fork}.

Hiện tượng đó được gọi là rẽ nhánh trạng thái mật mã (Cryptographic State Fork). Đây là tình huống trong đó các nút vẫn nghĩ rằng mình thuộc cùng một nhóm, thậm chí có thể cùng nhìn thấy cùng số epoch, nhưng thực tế lại đang nắm giữ các bí mật nhóm, hàm băm cây (tree hash) hoặc lịch sử bắt tay (transcript hash) khác nhau. Khi đó, thông điệp được tạo trên một nhánh không còn được xử lý hợp lệ trên nhánh còn lại. Khác với xung đột dữ liệu thông thường, trạng thái mật mã đã rẽ nhánh không thể được ghép lại bằng cách trộn hai phiên bản dữ liệu, vì mỗi nhánh đã sinh ra một chuỗi khóa và một lịch sử xác thực riêng \cite{rfc9420,alwen2023fork,kohbrok2025dmls}. Hiện tượng này được minh họa trong Hình \ref{fig:cryptographic_state_fork}.

\begin{figure}[htbp]
\centering
\resizebox{0.95\textwidth}{!}{%
\begin{tikzpicture}[
    font=\small,
    >=Latex,
    box/.style={
        draw,
        rectangle,
        rounded corners=3pt,
        align=center,
        minimum height=1.0cm,
        text width=4.2cm,
        inner sep=6pt
    },
    smallbox/.style={
        draw,
        rectangle,
        rounded corners=3pt,
        align=center,
        minimum height=0.9cm,
        text width=3.4cm,
        inner sep=5pt
    },
    note/.style={
        draw,
        rectangle,
        rounded corners=3pt,
        align=center,
        text width=11.0cm,
        inner sep=7pt
    },
    arrow/.style={->, thick}
]

\node[box] (epochE) {
    \textbf{Epoch $E$}\\
    Trạng thái nhóm chung $S_E$
};

\node[smallbox, below left=1.4cm and 1.8cm of epochE] (commitA) {
    Commit của nút $A$\\
    (dựa trên $S_E$)
};

\node[smallbox, below right=1.4cm and 1.8cm of epochE] (commitB) {
    Commit của nút $B$\\
    (dựa trên $S_E$)
};

\node[box, below=1.2cm of commitA] (branchA) {
    \textbf{Nhánh A}\\
    Epoch $E+1$\\
    Tree Hash: $TH_A$\\
    Transcript Hash: $H_A$
};

\node[box, below=1.2cm of commitB] (branchB) {
    \textbf{Nhánh B}\\
    Epoch $E+1$\\
    Tree Hash: $TH_B$\\
    Transcript Hash: $H_B$
};

\node[note, below=1.8cm of $(branchA)!0.5!(branchB)$] (result) {
    \textbf{Rẽ nhánh trạng thái mật mã}\\[1mm]
    Hai nhánh cùng mang số Epoch $E+1$ nhưng không còn cùng trạng thái MLS, vì vậy
    các nút ở hai nhánh không thể xử lý hợp lệ thông điệp của nhau.
};

\draw[arrow] (epochE.south) -- node[left, pos=0.45] {\scriptsize đồng thời} (commitA.north);
\draw[arrow] (epochE.south) -- node[right, pos=0.45] {\scriptsize đồng thời} (commitB.north);
\draw[arrow] (commitA.south) -- (branchA.north);
\draw[arrow] (commitB.south) -- (branchB.north);
\draw[arrow] (branchA.south) -- ++(0,-0.5) -| (result.north west);
\draw[arrow] (branchB.south) -- ++(0,-0.5) -| (result.north east);

\end{tikzpicture}
}
\caption{Minh họa hiện tượng rẽ nhánh trạng thái mật mã trong MLS}
\label{fig:cryptographic_state_fork}
\end{figure}

Trong các triển khai thông thường của MLS, việc hạn chế rẽ nhánh được hỗ trợ bởi Delivery Service (DS), có thể hiểu là lớp dịch vụ tiếp nhận và phân phối các thông điệp MLS giữa các thiết bị \cite{rfc9750}. Trong môi trường có DS nhất quán mạnh, khi nhiều thành viên cùng tạo Commit ở cùng một epoch, hệ thống có một điểm tự nhiên để quyết định Commit nào được xử lý trước. Nhờ đó, các thay đổi trạng thái dễ được quan sát theo một thứ tự thống nhất hơn.

Khó khăn xuất hiện khi đưa MLS sang mạng ngang hàng. Trong môi trường P2P, thông điệp có thể đến chậm, đến sai thứ tự, bị lặp lại hoặc tạm thời không đến được do thiết bị ngoại tuyến hay do mạng bị chia cắt. Khi đó, không còn một thực thể trung tâm nào tự nhiên đứng ra sắp xếp thứ tự cho các Commit. Nói cách khác, MLS đòi hỏi chuỗi epoch tiến hóa gần như tuyến tính, trong khi mạng ngang hàng bất đồng bộ lại không tự cung cấp thứ tự tuyến tính toàn cục cho các thông điệp thay đổi trạng thái \cite{rfc9750,gilbert2002brewers}. Đây chính là mâu thuẫn cốt lõi của bài toán mà đồ án cần giải quyết.

Trong phạm vi đồ án, hệ thống được xem như một mạng gồm nhiều thiết bị ngang hàng, mỗi thiết bị lưu cục bộ trạng thái MLS của các nhóm mà nó tham gia. Các tình huống được quan tâm gồm độ trễ mạng không đồng đều, thiết bị ngoại tuyến, sự biến động nút mạng và phân mảnh mạng. Mục tiêu của phần nghiên cứu không phải thay đổi các bảo đảm mật mã nền tảng của MLS, mà là tìm cách tổ chức việc điều phối ở tầng ứng dụng để các thiết bị vẫn có thể hội tụ trở lại một trạng thái nhóm hợp lệ sau những bất ổn của mạng.

\section{Các kết quả nghiên cứu tương tự}
\label{sec:nghien-cuu-lien-quan}

Các nghiên cứu liên quan có thể được nhìn từ ba hướng chính: các cơ chế nhắn tin nhóm xuất hiện trước MLS, MLS trong mô hình có Delivery Service và các nỗ lực đưa MLS sang môi trường phi tập trung. Việc phân chia theo hướng này giúp làm rõ vì sao đồ án không chọn lặp lại hoàn toàn một giải pháp có sẵn, mà cần một hướng tiếp cận riêng phù hợp hơn với mục tiêu P2P và Local-First.

\subsection{Các cơ chế mã hóa nhóm trước MLS}

Trước khi MLS được chuẩn hóa, nhiều hệ thống nhắn tin nhóm đã mở rộng từ cơ chế bảo mật hai bên sang môi trường nhiều người tham gia. Một nền tảng quan trọng là Double Ratchet, vốn rất hiệu quả cho liên lạc hai bên \cite{perrin2025doubleratchet}. Khi đưa sang nhóm, một hướng phổ biến là để mỗi người gửi tự tạo một khóa gửi tin riêng, rồi phân phối khóa đó cho các thành viên còn lại qua các kênh bảo mật hai bên. Cách làm này thường được gọi là Sender Keys \cite{balbas2023senderkeys}. Sau khi khóa gửi tin đã được chia sẻ, người gửi chỉ cần mã hóa một lần cho mỗi thông điệp, còn hạ tầng bên dưới chịu trách nhiệm phát tán bản mã tới cả nhóm. Ưu điểm của hướng này là chi phí gửi tin thấp và triển khai tương đối thực dụng. Tuy nhiên, khi nhóm đông hoặc thay đổi thành viên thường xuyên, chi phí cập nhật khóa và phục hồi sau thỏa hiệp vẫn tăng đáng kể.

Megolm trong hệ sinh thái Matrix cũng theo tinh thần tối ưu việc gửi tin nhóm \cite{ginesin2024matrix}. Có thể hiểu đơn giản rằng mỗi người gửi duy trì một phiên gửi ra (\textit{outbound session}), tức một trạng thái khóa được dùng liên tiếp để mã hóa nhiều tin nhắn trong cùng một phòng trò chuyện. Bên trong phiên đó, trạng thái khóa được cập nhật dần theo một cơ chế nhất định của nhóm (\textit{group ratchet}). Thiết kế này giúp việc gửi tin trong nhóm nhẹ hơn, nhưng điểm đổi lại là khả năng phục hồi sau thỏa hiệp không mạnh bằng những mô hình thay khóa nhóm chặt chẽ hơn. Nếu trạng thái của một phiên bị lộ, phạm vi thông điệp bị ảnh hưởng có thể kéo dài cho đến khi hệ thống tạo phiên mới và phân phối lại khóa.

Nhìn chung, Sender Keys và Megolm đều có giá trị thực tiễn cao đối với bài toán phát tán tin nhắn trong nhóm. Tuy nhiên, chúng không trực tiếp giải quyết bài toán mà đồ án quan tâm nhất, đó là duy trì một trạng thái nhóm nhất quán khi nhiều thiết bị có thể cùng thay đổi trạng thái trong môi trường P2P bất đồng bộ. Đây là lý do MLS trở thành nền tảng phù hợp hơn cho đề tài.

\subsection{MLS với Delivery Service}

MLS được chuẩn hóa để cung cấp cơ chế thiết lập khóa nhóm bất đồng bộ với các bảo đảm bảo mật mạnh hơn cho truyền thông nhóm quy mô lớn \cite{rfc9420}. Điểm quan trọng của MLS là nó không chỉ tối ưu việc gửi thông điệp, mà còn mô hình hóa rõ quá trình thay đổi thành viên và cập nhật khóa của cả nhóm thông qua các khái niệm như proposal, commit và epoch.

Trong kiến trúc chuẩn của MLS, Delivery Service thường đảm nhiệm việc lưu trữ một số vật liệu khóa công khai ban đầu, phân phối thông điệp MLS và hỗ trợ quá trình truyền thông giữa các client \cite{rfc9750}. Nhờ có một đầu mối như vậy, hệ thống triển khai theo mô hình Client-Server có điều kiện thuận lợi hơn để quan sát các Commit theo một thứ tự đủ nhất quán. Vì thế, MLS hoạt động tự nhiên hơn trong môi trường có DS, nhưng khi bỏ đi điểm điều phối này thì bài toán ordering của Commit lập tức trở nên khó hơn nhiều.

\subsection{Các hướng phi tập trung hóa MLS}

Một hướng tiếp cận trực tiếp là bổ sung một lớp đồng thuận phân tán để thống nhất thứ tự Commit. Các cơ chế như BFT, blockchain hoặc đồng thuận dựa trên quorum có thể tạo ra một nhật ký tuyến tính chung cho các thay đổi trạng thái nhóm. Nếu mọi Commit đều được ghi vào cùng một chuỗi thứ tự, nguy cơ rẽ nhánh trạng thái sẽ giảm rõ rệt. Tuy nhiên, đổi lại là chi phí truyền thông cao, phụ thuộc vào quorum và độ phức tạp triển khai lớn. Với mục tiêu xây dựng một hệ thống P2P nhẹ, có thể vận hành trong mạng nội bộ hoặc mạng không ổn định, đây không phải lựa chọn phù hợp \cite{gilbert2002brewers,localfirst2019}.

Một hướng khác là Decentralized MLS (DMLS) của Kohbrok, trong đó chính bản thân MLS được mở rộng để xử lý Commit đến ngoài thứ tự \cite{kohbrok2025dmls}. Về ý tưởng, client giữ lại thêm một phần vật liệu khóa hoặc trạng thái tạm thời để có thể xử lý những nhánh Commit khác nhau khi mạng phi tập trung không cung cấp ordering rõ ràng. Cách làm này cho thấy một hướng nghiên cứu đáng chú ý, nhưng đồng thời cũng làm tăng độ phức tạp triển khai và đặt ra thêm đánh đổi đối với Forward Secrecy \cite{alwen2023fork,kohbrok2025dmls}. Hơn nữa, hướng này hiện vẫn ở mức Internet-Draft, chưa trở thành chuẩn hoàn chỉnh.

Từ các hướng trên có thể thấy khoảng trống còn lại nằm ở chỗ sau: cần một cơ chế đủ nhẹ để phù hợp với P2P, nhưng vẫn đủ chặt để giảm nguy cơ rẽ nhánh và hỗ trợ hệ thống hội tụ trở lại sau khi mạng bị chia cắt. Đó cũng là động lực trực tiếp cho hướng tiếp cận mà đồ án theo đuổi ở các chương sau.

\section{Nền tảng MLS và TreeKEM}
\label{sec:nen-tang-mls-treekem}

\subsection{Các khái niệm cốt lõi: client, group, member và epoch}

Trong MLS, client là thực thể nắm giữ khóa mật mã và tham gia vào tiến trình thiết lập trạng thái nhóm. Group là một nhóm logic gồm các client cùng chia sẻ một trạng thái mật mã tại một thời điểm. Member là client đang thực sự nằm trong trạng thái chia sẻ đó và có quyền truy cập các bí mật của nhóm. Epoch, như đã trình bày ở trên, là một phiên bản cụ thể của trạng thái nhóm \cite{rfc9420}.

Điểm quan trọng cần nhấn mạnh là trạng thái nhóm trong MLS không phải một tập dữ liệu rời rạc có thể ghép lại tùy ý. Nó là một chuỗi epoch nối tiếp nhau, trong đó mỗi epoch mới được sinh ra từ epoch trước đó. Vì vậy, nếu các thành viên không cùng nhìn nhận một chuỗi epoch tương thích, nhóm sẽ mất khả năng duy trì một trạng thái mật mã chung.

\subsection{Ratchet Tree và TreeKEM}

TreeKEM là cơ chế trung tâm giúp MLS cập nhật khóa nhóm hiệu quả. Thay vì để mỗi thành viên phải trao đổi khóa riêng với tất cả các thành viên còn lại, MLS tổ chức các thành viên thành các lá của một cây nhị phân, gọi là Ratchet Tree. Có thể hiểu trực quan rằng mỗi lần nhóm thay đổi, hệ thống chỉ cần cập nhật một số bí mật dọc theo một đường trong cây, thay vì phát lại toàn bộ khóa cho cả nhóm \cite{rfc9420}.

Khi một thành viên cập nhật khóa, thành viên đó tạo ra một đường dẫn cập nhật (\textit{update path}) từ lá của mình lên gốc cây. Các bí mật mới trên đường đi này được mã hóa cho các nhánh đồng cấp tương ứng, thường được gọi là \textit{copath}. Nhờ cấu trúc cây, số lượng phép tính và lượng dữ liệu cập nhật chỉ tăng gần theo logarit của số thành viên trong nhóm. Nói ngắn gọn, TreeKEM là cơ chế giúp MLS thay khóa nhóm theo cách có cấu trúc và tiết kiệm hơn so với việc phân phối lại khóa theo kiểu tuyến tính cho toàn bộ thành viên.

Cần lưu ý rằng bí mật ở gốc cây không được dùng trực tiếp để mã hóa thông điệp ứng dụng. Sau mỗi Commit, MLS còn đưa các bí mật mới qua lịch trình sinh khóa (\textit{key schedule}) để tạo ra các khóa cụ thể cho từng mục đích như xác thực, mã hóa thông điệp hay xuất vật liệu khóa \cite{rfc9420}. Vì vậy, TreeKEM nên được hiểu là cơ chế cập nhật trạng thái khóa nhóm, còn việc sinh các khóa sử dụng trực tiếp cho thông điệp được xử lý ở bước sau.

\subsection{Proposal, Commit và tiến trình chuyển epoch}

MLS phân biệt rõ proposal và commit. Proposal là thông điệp đề xuất một thay đổi cho nhóm, chẳng hạn thêm thành viên, loại bỏ thành viên hoặc cập nhật khóa. Bản thân proposal chưa làm thay đổi epoch. Commit mới là thông điệp thực sự áp dụng một tập proposal và đưa nhóm sang epoch mới \cite{rfc9420}.

Nhìn dưới góc độ hệ phân tán, có thể xem Commit như thao tác ghi vào trạng thái mật mã của nhóm. Mỗi Commit đều làm xuất hiện một phiên bản trạng thái mới. Vì thế, nếu tại cùng một epoch mà có nhiều Commit cạnh tranh, nguy cơ rẽ nhánh trạng thái sẽ xuất hiện. Đây là lý do bài toán ordering của Commit luôn giữ vai trò trung tâm khi vận hành MLS trên mạng ngang hàng.

\subsection{Các dấu hiệu nhận biết trạng thái: tree hash, transcript hash và epoch authenticator}

Để kiểm tra các thành viên có còn cùng một trạng thái hay không, MLS sử dụng một số giá trị ràng buộc trạng thái \cite{rfc9420}. Trong đó, hàm băm cây (tree hash) phản ánh cấu trúc Ratchet Tree hiện tại; hàm băm lịch sử bắt tay (transcript hash) ràng buộc chuỗi thông điệp bắt tay đã được xác nhận; còn bộ xác thực epoch (epoch authenticator) có thể được dùng để đối chiếu việc các thành viên có thực sự đang chia sẻ cùng trạng thái epoch hay không.

Ý nghĩa của các giá trị này đối với đồ án là rất trực tiếp. Trong môi trường P2P, hai thiết bị có thể cùng báo cùng mã nhóm và cùng số epoch, nhưng điều đó chưa đủ để kết luận rằng chúng còn ở cùng trạng thái MLS. Nếu tree hash hoặc transcript hash khác nhau, về bản chất chúng đã tách thành hai nhánh. Vì vậy, khi nghiên cứu bài toán phát hiện phân nhánh và đồng bộ lại trạng thái, các dấu hiệu nhận biết này đóng vai trò đặc biệt quan trọng.

\subsection{Application messages, delayed messages và External Commit}

MLS phân biệt thông điệp bắt tay (\textit{handshake messages}) với thông điệp ứng dụng (\textit{application messages}) \cite{rfc9420}. Proposal và Commit thuộc nhóm thứ nhất vì chúng làm thay đổi hoặc chuẩn bị thay đổi trạng thái nhóm. Ngược lại, application messages là các tin nhắn ứng dụng được mã hóa bằng khóa sinh ra từ trạng thái của epoch hiện tại.

Trong mạng thực tế, thông điệp ứng dụng có thể đến muộn hoặc đến sai thứ tự. MLS có cơ chế hỗ trợ một mức độ nhất định cho các thông điệp bị trễ, chẳng hạn thông qua bộ khóa của người gửi (\textit{sender ratchet}) và bộ đếm thế hệ (\textit{generation counter}). Tuy nhiên, để bảo vệ Forward Secrecy, thiết bị không nên giữ các khóa cũ vô thời hạn. Điều này tạo ra một đánh đổi quen thuộc: nếu xóa khóa cũ quá sớm thì thông điệp đến muộn có thể không giải mã được; nếu giữ khóa cũ quá lâu thì phạm vi thiệt hại khi thiết bị bị thỏa hiệp sẽ tăng \cite{rfc9420}.

MLS cũng hỗ trợ cơ chế gia nhập từ bên ngoài, thường được gọi là External Commit hoặc External Join \cite{rfc9420}. Có thể hiểu đây là cách để một client chưa ở trong trạng thái hiện tại của nhóm gia nhập lại nhóm khi có đủ thông tin công khai cần thiết. Cơ chế này quan trọng đối với đồ án vì khi trạng thái đã rẽ nhánh, một thiết bị có thể cần từ bỏ nhánh cũ và tham gia lại nhánh đang được chọn.

\section{Nền tảng mạng ngang hàng và thứ tự trong hệ phân tán}
\label{sec:mang-ngang-hang-libp2p}

\subsection{Đặc tính của mạng ngang hàng}

Mạng ngang hàng khác với mô hình Client-Server ở chỗ không có một máy chủ trung tâm duy nhất điều phối toàn bộ luồng truyền thông. Các thiết bị có thể tự phát hiện nhau, thiết lập kết nối trực tiếp hoặc gián tiếp và trao đổi dữ liệu qua nhiều đường truyền khác nhau. Cách tổ chức này giúp hệ thống giảm phụ thuộc vào hạ tầng trung tâm, nhưng đồng thời làm cho việc điều phối trở nên khó hơn.

Trong môi trường P2P, các hiện tượng như thiết bị ngoại tuyến, độ trễ không đồng đều, trùng lặp thông điệp, biến động nút mạng hoặc phân mảnh mạng là điều bình thường. Với ứng dụng thông thường, chúng chủ yếu gây chậm hoặc mất dữ liệu tạm thời. Với MLS, các hiện tượng đó còn có thể kéo theo hệ quả nghiêm trọng hơn: một số thiết bị đã chuyển sang epoch mới trong khi thiết bị khác vẫn ở epoch cũ, hoặc hai phân vùng mạng cùng tạo ra những Commit khác nhau. Điều này cho thấy lớp mạng P2P chỉ giải quyết việc truyền thông, chứ không tự giải quyết bài toán nhất quán trạng thái mật mã.

\subsection{go-libp2p và GossipSub}

Trong nguyên mẫu của đồ án, tầng mạng được xây dựng trên go-libp2p \cite{libp2p2026doc}. Đây là thư viện P2P mô-đun hóa, cho phép mỗi thiết bị có một định danh mạng riêng, thiết lập kết nối qua nhiều cơ chế truyền dẫn và trao đổi dữ liệu trên nhiều luồng logic. Những khả năng này phù hợp với yêu cầu của một ứng dụng nhắn tin không phụ thuộc hoàn toàn vào máy chủ trung tâm.

Một thành phần quan trọng của libp2p là GossipSub, tức cơ chế phát tán thông điệp theo mô hình công bố/đăng ký nhận (\textit{publish/subscribe}). Có thể hiểu đơn giản rằng khi một thiết bị gửi thông điệp vào một chủ đề chung, thông điệp đó sẽ được lan truyền dần qua các thiết bị lân cận trong mạng. Cơ chế này phù hợp để phát tán proposal, commit hoặc các thông báo trạng thái trong môi trường P2P.

Tuy nhiên, GossipSub chỉ là cơ chế phát tán thông điệp, không phải cơ chế đồng thuận. Nó không bảo đảm mọi thiết bị nhận thông điệp theo cùng một thứ tự, không tự chọn Commit thắng cuộc khi có xung đột, và cũng không bảo đảm một thông điệp chỉ được nhận đúng một lần. Vì vậy, chỉ đưa MLS lên trên GossipSub là chưa đủ để bảo vệ chuỗi trạng thái của nhóm.

\subsection{Đồng hồ logic và thứ tự hiển thị thông điệp ứng dụng}

Trong cùng một epoch, nhiều thành viên có thể gửi thông điệp ứng dụng song song. Những thông điệp này không nhất thiết cần một thứ tự mật mã toàn cục như Commit, nhưng giao diện chat vẫn cần một thứ tự hiển thị đủ ổn định giữa các thiết bị. Nếu chỉ dựa vào đồng hồ vật lý, hệ thống dễ gặp sai lệch do lệch giờ hoặc do môi trường mạng không ổn định.

Vì vậy, hệ phân tán thường sử dụng các dạng đồng hồ logic để biểu diễn quan hệ thứ tự giữa các sự kiện. Lamport Clock cho một thứ tự logic đơn giản \cite{lamport1978}. Hybrid Logical Clock (HLC) kết hợp thời gian vật lý và bộ đếm logic để tạo ra dấu thời gian gần với thời gian thực hơn, nhưng vẫn ổn định hơn khi đồng hồ giữa các thiết bị không hoàn toàn khớp nhau \cite{kulkarni2014hlc}. Trong đồ án, loại dấu thời gian này chỉ có ý nghĩa đối với việc sắp xếp hiển thị các thông điệp ứng dụng; nó không thay thế epoch, chữ ký hay các cơ chế bảo vệ trạng thái của MLS.

\subsection{Khoảng trống cần giải quyết}

Từ các phân tích trên có thể thấy bài toán của đề tài nằm ở giao điểm giữa mật mã học nhóm và hệ phân tán. MLS cung cấp lõi mật mã mạnh nhưng giả định hệ thống triển khai phải giải quyết việc phân phối và ordering của thông điệp. Libp2p và GossipSub cung cấp hạ tầng truyền thông ngang hàng nhưng không cung cấp thứ tự tuyến tính cho Commit. HLC hỗ trợ sắp xếp thông điệp ứng dụng nhưng không bảo vệ trạng thái mật mã của nhóm \cite{rfc9750,libp2p2026doc,kulkarni2014hlc}.

Khoảng trống nghiên cứu vì vậy không nằm ở chỗ thiếu một giao thức mã hóa nhóm, mà nằm ở chỗ thiếu một cơ chế điều phối đủ nhẹ để phù hợp với P2P, nhưng vẫn đủ chặt để giảm nguy cơ rẽ nhánh và hỗ trợ hội tụ trở lại sau khi mạng bất ổn. Khoảng trống đó là cơ sở trực tiếp cho giải pháp sẽ được trình bày ở Chương 3.

Chương 2 đã làm rõ ngữ cảnh lý thuyết của bài toán vận hành MLS trên mạng ngang hàng. Trước hết, chương chỉ ra rằng khó khăn trung tâm của đề tài không chỉ là truyền tin trong môi trường P2P, mà là duy trì một trạng thái mật mã chung cho toàn nhóm khi không còn điểm tuần tự hóa trung tâm. Tiếp theo, chương khảo sát các hướng tiếp cận liên quan, từ Sender Keys, Megolm đến MLS trong mô hình có Delivery Service và các nỗ lực phi tập trung hóa MLS, qua đó làm nổi bật khoảng trống nghiên cứu mà đồ án hướng tới.

Bên cạnh đó, chương cũng trình bày các kiến thức nền tảng cần thiết để hiểu phần phương pháp ở các chương sau, bao gồm epoch, proposal, commit, TreeKEM, các dấu hiệu nhận biết trạng thái và cơ chế gia nhập lại nhóm. Cuối cùng, chương phân tích ngắn gọn các đặc tính của mạng ngang hàng, vai trò của libp2p, GossipSub và đồng hồ logic trong việc hỗ trợ truyền thông và sắp xếp thông điệp ứng dụng. Những nội dung này tạo tiền đề trực tiếp cho Chương 3, nơi giải pháp điều phối phi tập trung của đồ án sẽ được trình bày chi tiết.

\end{document}
